\documentclass[12pt]{article}

\usepackage{fontspec}
\newfontfamily\handwriting{a-casual-handwritten-pen}[Path=../assets/fonts/,Extension=.ttf,Scale=1.5]
\newfontfamily\comic{Comic Neue}

\usepackage[utf8]{inputenc}
\usepackage{amsmath}
\usepackage{amssymb}
\usepackage{enumitem}
\usepackage{xcolor}
\usepackage{soul}
\usepackage{tikz}
\usepackage[most]{tcolorbox}
\usepackage{titlesec}
\usepackage{multicol}
\usepackage{environ}

\usetikzlibrary{fit, positioning, arrows.meta, decorations.pathreplacing}

% --- Custom Colors ---
\definecolor{primaryblue}{RGB}{42, 111, 151}
\definecolor{exbg}{RGB}{255, 240, 245}
\definecolor{rulebg}{RGB}{232, 245, 233}
\definecolor{highlight}{RGB}{255, 243, 176}
\definecolor{pinkanno}{RGB}{255, 105, 180}
\definecolor{purpleanno}{RGB}{147, 112, 219}
\definecolor{solutionbg}{RGB}{245, 245, 255}

% --- Header Styling ---
\titleformat{\section}{\color{primaryblue}\normalfont\Large\bfseries}{\thesection}{1em}{}
\titleformat{\subsection}{\color{primaryblue}\normalfont\large\bfseries}{\thesubsection}{1em}{}

% --- Friendly Boxes ---
% --- Auto-numbered Example Box ---
\newcounter{examplenum}

\newtcolorbox{friendlyexbox}[1]{
    colback=exbg,
    colframe=pinkanno!50,
    arc=5pt,
    title=\textbf{#1},
    coltitle=black,
    attach title to upper,
    after title={:\quad}
}

\newenvironment{friendlyex}{%
    \refstepcounter{examplenum}%
    \begin{friendlyexbox}{Example \theexamplenum}%
}{%
    \end{friendlyexbox}%
}

\newtcolorbox{friendlyrule}[1]{
    colback=rulebg,
    colframe=primaryblue!30,
    arc=5pt,
    fonttitle=\bfseries,
    title={#1},
    coltitle=primaryblue
}

\newcommand{\funhl}[1]{\sethlcolor{highlight}\hl{#1}}

% --- Show/Hide Solutions ---
\newif\ifshowsolutions

% Uncomment the next line to show solutions.
\showsolutionstrue

\newsavebox{\solutionmeasurebox}

\NewEnviron{solutionbox}{
\ifshowsolutions
\begin{tcolorbox}[
    colback=solutionbg,
    colframe=purpleanno!45,
    arc=5pt,
    title={Solution},
    coltitle=primaryblue,
    fonttitle=\bfseries
]
\BODY
\end{tcolorbox}
\else
\sbox{\solutionmeasurebox}{\begin{minipage}{\dimexpr\linewidth-10pt\relax}\BODY\end{minipage}}
\vspace{\dimexpr\ht\solutionmeasurebox+\dp\solutionmeasurebox+3em\relax}
\fi
}

\begin{document}
\handwriting

\begin{center}
    \Large\textbf{\color{primaryblue}Module 4 Lesson 4A\\Logarithmic Equations and Applications} \\
    \vspace{0.2cm}
    \hrule
\end{center}

\section*{Objectives}

\begin{friendlyrule}{In this lesson, we will learn how to}
\begin{itemize}
    \item solve logarithmic equations using the equivalence property;
    \item combine logarithms first when an equation has multiple log terms;
    \item check solutions against the domain of the original logarithms;
    \item apply logarithmic equations to the pH scale.
\end{itemize}
\end{friendlyrule}

\section*{1. Equivalence Property}

\begin{friendlyrule}{Equivalence Property of Logarithmic Equations}
Let $a,x,y$ be positive real numbers with $a\neq1$. If
\[
\log_a x=\log_a y \qquad \text{then} \qquad x=y.
\]
\end{friendlyrule}

\begin{friendlyrule}{A Reminder}
Since $\log_a(x)$ is only defined for $x>0$, always check that a candidate solution keeps every logarithm's argument \textbf{positive} in the original equation. A solution that fails this check is \textbf{extraneous} and must be discarded.
\end{friendlyrule}

\begin{friendlyex}{}
Solve.
\begin{enumerate}[label=(\alph*)]
    \item $\log(x^2)=\log(16)$
    \item $\log_5(8-7x)=3$
\end{enumerate}
\end{friendlyex}

\begin{solutionbox}
\textbf{(a)} By the equivalence property, $x^2=16 \implies x=\pm4$. Both keep $x^2=16>0$, so both are valid.
\[
\boxed{x=4,\ x=-4}
\]

\textbf{(b)} Rewrite in exponential form: $8-7x=5^3=125$.
\[
-7x=117 \implies x=-\frac{117}{7}
\]
Check: $8-7x=125>0$. \checkmark
\[
\boxed{x=-\frac{117}{7}}
\]
\end{solutionbox}

\section*{2. Combining Logarithms Before Solving}

When an equation has multiple logarithmic terms, first combine them into a single logarithm using the product/quotient/power rules, then apply the equivalence property.

\begin{friendlyex}{}
Solve $\ln x-\ln(x-4)=\ln 3$.
\end{friendlyex}

\begin{solutionbox}
\[
\ln\left(\frac{x}{x-4}\right)=\ln3 \implies \frac{x}{x-4}=3 \implies x=3(x-4)=3x-12.
\]
\[
-2x=-12 \implies x=6.
\]

Check: $x=6>0$ and $x-4=2>0$. \checkmark
\[
\boxed{x=6}
\]
\end{solutionbox}

\begin{friendlyex}{}
Solve $\log_2(x-4)=\log_2(2)-\log_2(x-3)$.
\end{friendlyex}

\begin{solutionbox}
Combine the right side: $\log_2(2)-\log_2(x-3)=\log_2\left(\dfrac{2}{x-3}\right)$.
\[
\log_2(x-4)=\log_2\left(\frac{2}{x-3}\right) \implies x-4=\frac{2}{x-3} \implies (x-4)(x-3)=2
\]
\[
x^2-7x+12=2 \implies x^2-7x+10=0 \implies (x-5)(x-2)=0.
\]
\[
x=5 \text{ or } x=2.
\]

Check $x=5$: $x-4=1>0$ and $x-3=2>0$. \checkmark

Check $x=2$: $x-4=-2<0$. \textbf{Fails} --- extraneous.
\[
\boxed{x=5}
\]
\end{solutionbox}

\begin{friendlyex}{}
Solve $2\ln x-\ln 5=\ln(x+10)$.
\end{friendlyex}

\begin{solutionbox}
\[
\ln(x^2)-\ln5=\ln(x+10) \implies \ln\left(\frac{x^2}{5}\right)=\ln(x+10) \implies \frac{x^2}{5}=x+10
\]
\[
x^2=5x+50 \implies x^2-5x-50=0 \implies (x-10)(x+5)=0.
\]
\[
x=10 \text{ or } x=-5.
\]

Check $x=10$: $x>0$ and $x+10=20>0$. \checkmark

Check $x=-5$: $\ln x$ requires $x>0$, but $x=-5<0$. \textbf{Fails} --- extraneous.
\[
\boxed{x=10}
\]
\end{solutionbox}

\begin{friendlyex}{}
Solve $\log_3(\log_2 x)=0$.
\end{friendlyex}

\begin{solutionbox}
Rewrite in exponential form (base $3$): $\log_2 x=3^0=1$.

Rewrite again (base $2$): $x=2^1=2$.
\[
\boxed{x=2}
\]
\end{solutionbox}

\section*{3. Application: The pH Scale}

\begin{friendlyrule}{pH Formula}
The acidity of a solution is measured using the pH scale, defined by
\[
pH=-\log[H^+],
\]
where $[H^+]$ is the concentration of the hydrogen ion in moles per liter.
\end{friendlyrule}

\begin{friendlyex}{}
A certain solution has pH of $3.5$. What is the hydrogen ion concentration $[H^+]$ in the solution?
\end{friendlyex}

\begin{solutionbox}
\[
3.5=-\log[H^+] \implies \log[H^+]=-3.5 \implies [H^+]=10^{-3.5}
\]
\[
\boxed{[H^+]\approx 3.16\times10^{-4} \text{ moles/liter}}
\]
\end{solutionbox}

\section*{Summary}

\begin{friendlyrule}{Key Ideas}
\begin{itemize}
    \item Equivalence Property: if $\log_a x=\log_a y$ (same base), then $x=y$.
    \item If an equation has a single logarithm equal to a number, rewrite it in exponential form: $\log_a(x)=p \implies x=a^p$.
    \item If an equation has multiple logarithmic terms, first combine them into a single logarithm with the product/quotient/power rules, \textbf{then} apply the equivalence property.
    \item \textbf{Always check every solution} against the domain of each original logarithm ($\text{argument}>0$) --- solving often introduces extraneous roots that must be discarded.
    \item Application: $pH=-\log[H^+]$, so $[H^+]=10^{-pH}$.
\end{itemize}
\end{friendlyrule}

\end{document}
